\documentclass[11pt, a4paper]{article}

\usepackage[utf8]{inputenc}
\usepackage[spanish]{babel}
\usepackage{amsmath, amssymb}
\usepackage{geometry}
\usepackage{fancyhdr}

\geometry{left=2cm, right=2cm, top=2.5cm, bottom=2.5cm}
\setlength{\headheight}{14pt}

\pagestyle{fancy}
\fancyhf{}
\lhead{\textbf{Universidad Católica Boliviana ``San Pablo''}}
\rhead{\textbf{IMT-121: Circuitos Electrónicos I}}
\cfoot{\thepage}

\begin{document}

\begin{center}
    \Large \textbf{Evaluación Diagnóstica de Competencias Previas} \\
    \vspace{0.2cm}
    \normalsize \textit{Semestre II/2026 -- Evaluación Formativa (0\% Ponderación Sumativa)}
\end{center}

\vspace{0.5cm}
\noindent
\textbf{Nombre y Apellido:} \underline{\hspace{7.5cm}} \hfill \textbf{Fecha:} 03/08/2026 \\
\textbf{Programa Académico:} $\square$ Ing. Mecatrónica \quad $\square$ Ing. Biomédica \hfill \textbf{Firma:} \underline{\hspace{2.5cm}}

\vspace{0.3cm}
\hrule
\vspace{0.5cm}

\noindent \textit{\textbf{Directriz Académica:} El propósito de este instrumento es auditar el dominio de las competencias previas (MAT-123 Álgebra Lineal). Desarrolle las soluciones de manera analítica; el uso de calculadoras con resolución matricial integrada está estrictamente restringido.}

\vspace{0.5cm}

\section*{Sección I: Álgebra Lineal Aplicada}

\textbf{Problema 1.} El siguiente sistema de ecuaciones lineales describe la distribución de corrientes ($I_1, I_2, I_3$) en los lazos de una red eléctrica resistiva estacionaria. Exprese dicho sistema en su formulación matricial canónica $\mathbf{A}\mathbf{x} = \mathbf{b}$.

\begin{align*}
    12 I_1 - 4 I_2 - 8 I_3 &= 10 \\
    -4 I_1 + 10 I_2 - 2 I_3 &= 0 \\
    -8 I_1 - 2 I_2 + 15 I_3 &= -5
\end{align*}

\vspace{2.5cm}

\textbf{Problema 2.} Refiriéndose a la matriz de coeficientes $\mathbf{A}$ obtenida en el inciso anterior:
\begin{enumerate}
    \item Defina la condición matemática fundamental para que el sistema admita una solución única.
    \item Analice la implicación matemática sobre el sistema si el determinante de $\mathbf{A}$ resultara ser estrictamente igual a cero ($\det(\mathbf{A}) = 0$).
\end{enumerate}

\vspace{3.5cm}

\section*{Sección II: Fenomenología Eléctrica Básica}

\textbf{Problema 3.} Establezca la distinción física formal entre la magnitud de \textit{corriente eléctrica} y la magnitud de \textit{diferencia de potencial} (voltaje).

\vspace{3.5cm}

\textbf{Problema 4.} Considere el siguiente escenario de instrumentación: Un operario configura un multímetro digital en función amperimétrica para medir corriente. Acto seguido, conecta las sondas de prueba \textbf{en paralelo} directamente a los terminales de una fuente de tensión ideal de $12\text{ V}$. 
\\
Describa el fenómeno físico resultante en el instrumento, fundamentando su respuesta matemáticamente a través de la Ley de Ohm.

\end{document}